\let\negmedspace\undefined
\let\negthickspace\undefined
\documentclass[journal]{IEEEtran}
\usepackage[a5paper, margin=10mm, onecolumn]{geometry}
%\usepackage{lmodern} % Ensure lmodern is loaded for pdflatex
\usepackage{tfrupee} % Include tfrupee package

\setlength{\headheight}{1cm} % Set the height of the header box
\setlength{\headsep}{0mm}     % Set the distance between the header box and the top of the text

\usepackage{gvv-book}
\usepackage{gvv}
\usepackage{cite}
\usepackage{amsmath,amssymb,amsfonts,amsthm}
\usepackage{algorithmic}
\usepackage{graphicx}
\usepackage{textcomp}
\usepackage{xcolor}
\usepackage{txfonts}
\usepackage{listings}
\usepackage{enumitem}
\usepackage{mathtools}
\usepackage{gensymb}
\usepackage{comment}
\usepackage[breaklinks=true]{hyperref}
\usepackage{tkz-euclide} 
\usepackage{listings}
% \usepackage{gvv}                                        
\def\inputGnumericTable{}                                 
\usepackage[latin1]{inputenc}                                
\usepackage{color}                                            
\usepackage{array}                                            
\usepackage{longtable}                                       
\usepackage{calc}                                             
\usepackage{multirow}                                         
\usepackage{hhline}                                           
\usepackage{ifthen}                                           
\usepackage{lscape}
\begin{document}

\bibliographystyle{IEEEtran}

\title{9.6.3}
\author{EE25BTECH11054 - S.Harsha Vardhan}
% \maketitle
% \newpage
% \bigskip
\maketitle 
\renewcommand{\thefigure}{\theenumi}
\renewcommand{\thetable}{\theenumi}
\setlength{\intextsep}{10pt} % Space between text and floats

\numberwithin{align}{enumi}
\numberwithin{figure}{enumi}
\renewcommand{\thetable}{\theenumi}
\vspace{-1em}
\textbf{Question 9.6.3}
Find the area of the circle $x^2 + y^2 = 16$ exterior to the parabola $y^2 = 6x$.

\solution
The total area of the circle $x^2 + y^2 = 16$ is $A_{circle} = \pi r^2 = 16\pi$.
First, we find the points of intersection of the two curves.
The conic parameters for the two curves are:
For the circle $x^2 + y^2 - 16 = 0$:
\begin{align}
\mtx{V_1} = \myvec{1 & 0 \\ 0 & 1}, \quad \vec{u_1} = \myvec{0 \\ 0}, \quad f_1 = -16. \tag{1}
\end{align}
For the parabola $y^2 - 6x = 0$:
\begin{align}
\mtx{V_2} = \myvec{0 & 0 \\ 0 & 1}, \quad \vec{u_2} = \myvec{-3 \\ 0}, \quad f_2 = 0. \tag{2}
\end{align}
The intersection of the two conics is given by the equation:
\begin{equation}
	\vec{x}^\top(\mtx{V_1} + \mu\mtx{V_2})\vec{x} + 2(\vec{u_1} + \mu\vec{u_2})^\top\vec{x} + (f_1 + \mu f_2) = 0 \tag{3}
\end{equation}
For a degenerate conic, the determinant of the quadratic part's matrix must be zero.
\begin{align}
\det(\mtx{V_1} + \mu\mtx{V_2}) &= 0 \tag{4} \\
\det\left(\myvec{1 & 0 \\ 0 & 1} + \mu\myvec{0 & 0 \\ 0 & 1}\right) &=
\begin{vmatrix}
	1 & 0 \\
	0 & 1+\mu
\end{vmatrix}
 = 1(1+\mu) = 0 \tag{5} \\
\implies \mu &= -1 \tag{6}
\end{align}
Substituting $\mu = -1$ into equation (3):
\begin{align} 
\vec{x}^\top( \mtx{V_1} - \mtx{V_2} )\vec{x} + 2( \vec{u_1} - \vec{u_2} )^\top\vec{x} + ( f_1 - f_2 ) &= 0 \tag{7} \\
\vec{x}^\top\left( \myvec{1 & 0 \\ 0 & 1} - \myvec{0 & 0 \\ 0 & 1} \right)\vec{x} + 2\left( \myvec{0 \\ 0} - \myvec{-3 \\ 0} \right)^\top\vec{x} + ( -16 - 0 ) &= 0 \tag{8} \\
\vec{x}^\top \myvec{1 & 0 \\ 0 & 0} \vec{x} + 2 \myvec{3 \\ 0}^\top \vec{x} - 16 &= 0 \tag{9}
\end{align}
Letting $\vec{x} = \myvec{x \\ y}$, the equation becomes:
\begin{align}
\myvec{x & y} \myvec{1 & 0 \\ 0 & 0} \myvec{x \\ y} + 2 \myvec{3 & 0} \myvec{x \\ y} - 16 = 0 \tag{10}
\end{align}
\begin{align}
\implies x^2 + 6x - 16 = 0 \tag{11}
\end{align}
Solving for $x$: $(x+8)(x-2) = 0$, which gives $x=2$ or $x=-8$.
Substituting these values into the parabola's equation $y^2 = 6x$:
If $x=2$, $y^2 = 12 \implies y = \pm\sqrt{12} = \pm 2\sqrt{3}$.
If $x=-8$, $y^2 = -48$, which has no real solution for $y$.
Hence, the points of intersection are:
\begin{align} 
\vec{a_1} = \myvec{2 \\ 2\sqrt{3}}, \quad \vec{a_2} = \myvec{2 \\ -2\sqrt{3}} \tag{12}
\end{align}
\[ A_{interior} = 2 \left[ \int_{0}^{2} \sqrt{6x} \,dx + \int_{2}^{4} \sqrt{16 - x^2} \,dx \right] \tag{13} \]
First integral:
\begin{align}
\int_{0}^{2} \sqrt{6x} \,dx = \sqrt{6} \left[ \frac{2}{3}x^{3/2} \right]_0^2 = \frac{2\sqrt{6}}{3}(2\sqrt{2}) = \frac{4\sqrt{12}}{3} = \frac{8\sqrt{3}}{3} \tag{14}
\end{align}
Second integral (using $\int \sqrt{a^2-x^2} \,dx = \frac{x}{2}\sqrt{a^2-x^2} + \frac{a^2}{2}\sin^{-1}\frac{x}{a}$):
\begin{align}
\int_{2}^{4} \sqrt{16-x^2} \,dx &= \left[ \frac{x}{2}\sqrt{16-x^2} + \frac{16}{2}\sin^{-1}\left(\frac{x}{4}\right) \right]_2^4 \nonumber \\
&= \left( 0 + 8\sin^{-1}(1) \right) - \left( \frac{2}{2}\sqrt{16-4} + 8\sin^{-1}\left(\frac{1}{2}\right) \right) \nonumber \\
&= 8\left(\frac{\pi}{2}\right) - \left(\sqrt{12} + 8\left(\frac{\pi}{6}\right)\right) = 4\pi - 2\sqrt{3} - \frac{4\pi}{3} = \frac{8\pi}{3} - 2\sqrt{3} \tag{15}
\end{align}
Adding the results for $A_{interior}$:
\begin{align}
A_{interior} &= 2 \left[ \frac{8\sqrt{3}}{3} + \frac{8\pi}{3} - 2\sqrt{3} \right] = 2 \left[ \frac{2\sqrt{3}}{3} + \frac{8\pi}{3} \right] = \frac{4\sqrt{3}}{3} + \frac{16\pi}{3} \tag{16}
\end{align}
The required area exterior to the parabola is $A_{exterior} = A_{circle} - A_{interior}$.
\begin{align}
A_{exterior} = 16\pi - \left( \frac{4\sqrt{3}}{3} + \frac{16\pi}{3} \right) = \frac{48\pi - 16\pi}{3} - \frac{4\sqrt{3}}{3} = \frac{32\pi - 4\sqrt{3}}{3} \tag{17}
\end{align}
\begin{figure}[h]
    \centering
    \includegraphics[width=1.1\linewidth]{figs/Figure_6.png}
    \caption{}
    \label{fig:placeholder}
\end{figure}

\end{document}
